\chapter*{Glossário}

\begin{description}

    \item[Edge Computing] É uma framework distribuída que aproxima as aplicações
    corporativas das fontes de dados, como os dispositivos de IoT ou
    servidores de \textit{edge} locais.

    \item[IoT (\textit{Internet of Things})] Refere-se à rede de objetos físicos, como sensores, câmaras e atuadores, que contêm tecnologia incorporada para comunicar e interagir através da Internet.

    \item[Machine Learning] Aprendizagem Automática. É um subcampo da Inteligência Artificial que utiliza algoritmos e modelos estatísticos para que os computadores aprendam a executar tarefas a partir de dados.

    \item[LightGBM] \textit{Light Gradient Boosting Machine}. É um algoritmo de \textit{Machine Learning} baseado em árvores de decisão, caracterizado pela sua elevada velocidade de treino, baixo consumo de memória e adequação para \textit{Edge Computing}.

    \item[XGBoost] \textit{Extreme Gradient Boosting}. É um algoritmo de \textit{Machine Learning} altamente eficiente e preciso, amplamente utilizado em competições de ciência de dados para problemas de classificação.

    \item[MQTT] \textit{Message Queuing Telemetry Transport}. É um protocolo de mensagens leve, baseado no modelo \textit{Publisher-Subscriber}, concebido especificamente para dispositivos com recursos limitados e redes de baixa largura de banda.

    \item[GNS3] \textit{Graphical Network Simulator-3}. É um emulador de redes de computadores que permite simular topologias complexas recorrendo a máquinas virtuais e contentores reais em vez de simulações simplificadas.

    \item[Docker] Plataforma de \textit{software} que permite empacotar, distribuir e executar aplicações dentro de contentores isolados. Consome muito menos recursos (RAM e Disco) em comparação com Máquinas Virtuais tradicionais.

    \item[iptables] Utilitário de linha de comandos do sistema operativo Linux que permite a um administrador de sistemas configurar as regras do \textit{firewall} fornecidas pelo \textit{kernel}, controlando o tráfego de rede entrante e sainte.

    \item[DDoS] \textit{Distributed Denial of Service} (Ataque Distribuído de Negação de Serviço). É um ataque malicioso em que múltiplos sistemas comprometidos são usados para inundar um alvo com tráfego, tornando-o indisponível para utilizadores legítimos.

    \item[Mirai Botnet] É um \textit{malware} notório que infeta sistemas baseados em Linux, especialmente dispositivos IoT antigos ou com credenciais de fábrica, transformando-os numa rede de \textit{bots} (Zombies) para lançar ataques DDoS em larga escala.

    \item[Spoofing] Técnica de falsificação onde um atacante se faz passar por outro dispositivo na rede (falsificando o endereço IP ou MAC) para obter acesso não autorizado ou desviar tráfego.

    \item[Zero-Day] Refere-se a ameaças ou anomalias que o sistema de segurança ou o modelo de Inteligência Artificial nunca viu durante a sua fase de treino, testando a sua capacidade de generalização e deteção de ataques inéditos.

    \item[TCP/IP] Conjunto de protocolos de comunicação fundamentais da Internet. Resulta da união entre o protocolo de controlo de transmissão (TCP) e o protocolo de inter-rede (IP).

    \item[Firewall] Dispositivo físico ou \textit{software} (como o \textit{iptables}) desenvolvido para aplicar políticas de segurança, bloqueando ou permitindo o tráfego de dados com base num conjunto de regras predefinidas.

\end{description}